% TOP_PROBLEM: {"id": 8, "title": "Odd Perfect Numbers", "category_id": 1, "category": {"id": 1, "name": "number_theory", "display_name": "Number Theory", "description": "Properties of integers, prime numbers, Diophantine equations.", "slug": "number-theory", "order_index": 1, "created_at": "2026-07-31T15:26:25.670Z"}, "status": "open"}
% FIRST_POSTED: 2026-09-25
% ENTRY_KIND: statement_only
% !TeX program = lualatex
% Definition and sources only; this file is not an LLM solution attempt.
\documentclass[11pt]{article}
\usepackage[margin=25mm]{geometry}
\usepackage{fontspec}
\setmainfont{Latin Modern Roman}
\usepackage{amsmath,amssymb,mathtools}
\usepackage{unicode-math}
\setmathfont{Latin Modern Math}
\usepackage{xurl,hyperref}
\hypersetup{colorlinks=true,urlcolor=blue}
\setcounter{secnumdepth}{0}
\setlength{\parindent}{0pt}
\setlength{\parskip}{5pt}
\setlength{\emergencystretch}{3em}
\begin{document}
\section*{\texorpdfstring{Odd Perfect Numbers}{Odd Perfect Numbers}}\label{8}
\subsection{Definitions and mathematical statement}
For an integer $N\ge1$, let $\sigma(N)=\sum_{d\mid N,\ d>0}d$. A perfect number satisfies $\sigma(N)=2N$, equivalently the sum of its positive proper divisors equals $N$. The question is
\[
 \exists N\in{\mathbb{Z}}_{>0}\quad N\equiv1\pmod2\quad\text{and}\quad\sigma(N)=2N\ ?
\]
For a prime factorization $N=\prod p_i^{a_i}$, the condition is
$\prod_i(1+p_i+\cdots+p_i^{a_i})=2\prod_i p_i^{a_i}$, with all $p_i$ odd. Either an example or an impossibility proof resolves the existence question; excluding any finite size range alone does not.
\par\smallskip\textit{Formulation and concept references:} \href{https://www.ams.org/mcom/2015-84-295/S0025-5718-2015-02941-X/S0025-5718-2015-02941-X.pdf}{[S1]}.

\subsection{Short English statement}
Can an odd positive integer equal the sum of all its positive proper divisors?
\subsection{Sources}
\begin{itemize}
\item[S1]Pace P. Nielsen, Mathematics of Computation 84 (2015), 2549-2567.
\url{https://www.ams.org/mcom/2015-84-295/S0025-5718-2015-02941-X/S0025-5718-2015-02941-X.pdf}.
\textit{Formulation source.}
\end{itemize}
\end{document}
